% Part: first-order-logic
% Chapter: lindstrom
% Section: abstract-logics

\documentclass[../../../include/open-logic-section]{subfiles}

\begin{document}

\olfileid[pt]{mod}{lin}{alg}
\olsection{Lógicas abstratas}

\begin{defn}
Uma \emph{lógica abstrata} é um par $\tuple{L, \models_L}$, onde $L$ é
uma função que atribui a cada !!{language}~$\Lang{L}$ um conjunto
$L(\Lang{L})$ de !!{sentence}s, e $\models_L$ é uma relação entre
!!{structure}s para a !!{language}~$\Lang{L}$ e !!{element}s de
$L(\Lang{L})$. Em particular, $\tuple{F, \models}$ é a lógica usual de
primeira ordem, isto é, $F$ é a função que atribui à
!!{language}~$\Lang{L}$ o conjunto das !!{sentence}s de primeira ordem
construídas a partir das constantes em $\Lang{L}$, e $\models$ é a
relação de satisfação da lógica de primeira ordem.
\end{defn}

Observe que ainda empregamos a mesma noção de !!{structure} para uma
!!{language} dada que na lógica de primeira ordem, mas não pressupomos
que as !!{sentence}s sejam construídas da maneira usual a partir dos
símbolos básicos em $\Lang{L}$, nem que a relação $\models_L$ seja
definida recursivamente como na lógica de primeira ordem. Assim, por
exemplo, por ser inteiramente geral, a definição pretende abranger o
caso em que as !!{sentence}s em $\tuple{L,\models_L}$ contêm conjunções
ou disjunções infinitamente longas, quantificadores diferentes de
$\lexists$ e $\lforall$ (por exemplo, ``existem infinitos~$x$ tais
que \dots''), ou talvez prefixos de quantificadores infinitamente
longos. Para enfatizar que as ``!!{sentence}s'' em $L(\Lang{L})$ não
precisam ser !!{sentence}s usuais da lógica de primeira ordem, neste
capítulo usamos as !!{variable}s $!E$, $!F$,~\dots para percorrê-las e
reservamos $!A$, $!B$,~\dots para !!{formula}s usuais de primeira ordem.

\begin{defn}
Seja $\Mod(L){!E}$ a classe $\Setabs{\Struct{M}}{\Struct{M}
  \models_L !E}$. Se for preciso explicitar a !!{language}, escrevemos
$\Mod[L](L){!E}$. Duas !!{structure}s $\Struct{M}$ e $\Struct{N}$ para
$\Lang{L}$ são \emph{elementarmente equivalentes em}
$\tuple{L, \models_L}$, o que escrevemos
$\Struct{M} \elemequiv[L] \Struct{N}$, se as mesmas !!{sentence}s de
$L(\Lang{L})$ são verdadeiras em ambas.
\end{defn}

\begin{defn}
Uma lógica abstrata $\tuple{L,\models_L}$ para a !!{language}
$\Lang{L}$ é \emph{normal} se satisfaz as seguintes propriedades:
\begin{enumerate}
\item (\emph{Monotonicidade de $L$}) Para as !!{language}s $\Lang{L}$
  e $\Lang{L'}$, se $\Lang{L} \subseteq \Lang{L'}$, então
  $L(\Lang{L}) \subseteq L(\Lang{L'})$.
\item (\emph{Propriedade da expansão}) Para cada $!E \in L(\Lang{L})$,
  existe um subconjunto \emph{finito} $\Lang{L'}$ de $\Lang{L}$ tal
  que a relação $\Struct{M} \models_L !E$ depende apenas do reduto de
  $\Struct{M}$ a $\Lang{L'}$; isto é, se $\Struct{M}$ e $\Struct{N}$
  têm o mesmo reduto a $\Lang{L'}$, então
  $\Struct{M} \models_L !E$ se, e somente se,
  $\Struct{N} \models_L !E$.
\item (\emph{Propriedade do isomorfismo}) Se
  $\Struct{M} \models_L !E$ e $\Struct{M} \simeq \Struct{N}$, então
  também $\Struct{N} \models_L !E$.
\item (\emph{Propriedade da renomeação}) A relação $\models_L$ é
  preservada sob renomeação: se a !!{language} $\Lang{L}'$ é obtida de
  $\Lang{L}$ substituindo cada símbolo $P$ por um símbolo $P'$ da mesma
  aridade e cada constante $c$ por uma constante distinta $c'$, então,
  para cada !!{structure}~$\Struct{M}$ e !!{sentence}~$!E$,
  $\Struct{M} \models_L !E$ se, e somente se,
  $\Struct{M}' \models_L !E'$, onde $\Struct{M}'$ é a
  $\Lang{L}'$-!!{structure} correspondente a $\Lang{L}$ e
  $!E' \in L(\Lang{L}')$.
\item (\emph{Propriedade booleana}) A lógica abstrata
  $\tuple{L, \models_L}$ é fechada sob os conectivos booleanos no
  seguinte sentido: para cada $!E \in L(\Lang{L})$ existe
  $!F \in L(\Lang{L})$ tal que $\Struct{M} \models_L !F$ se, e somente
  se, $\Struct{M} \not\models_L !E$; e, para cada $!E$ e $!F$, existe
  $!G$ tal que $\Mod(L){!G} = \Mod(L){!E} \cap \Mod(L){!F}$. O mesmo
  vale para !!{formula}s atômicas e para os demais conectivos.
\item (\emph{Propriedade dos quantificadores}) Para cada constante $c$
  em $\Lang{L}$ e cada $!E \in L(\Lang{L})$, existe
  $!F \in L(\Lang{L})$ tal que
  \[
  \Mod[L'](L){!F} = \Setabs{\Struct{M}}{\Expan{M}{a} \in
  \Mod[L](L){!E} \text{ para algum } a \in \Domain{M}},
  \]
  onde $\Lang{L'} = \Lang{L} \setminus \{c\}$ e $\Expan{M}{a}$ é a
  expansão de $\Struct{M}$ a $\Lang{L}$ que atribui $a$ a~$c$.
\item (\emph{Propriedade da relativização}) Dada !!a{sentence}
  $!E \in L(\Lang{L})$ e símbolos $R$, $c_1$, \dots, $c_n$ que não
  pertencem a $\Lang{L}$, existe !!a{sentence}
  $!F \in L(\Lang{L} \cup \{R,c_1,\ldots,c_n\})$, chamada
  \emph{relativização} de $!E$ a $\Atom{R}{x, c_1, \dots c_n}$, tal
  que, para cada !!{structure}~$\Struct{M}$:
  \[
  \Expan{M}{X, b_1, \ldots, b_n} \models_L !F \text{ se, e
    somente se, } \Struct{N} \models_L !E,
  \]
  onde $\Struct{N}$ é a subestrutura de $\Struct{M}$ com !!{domain}
  $\Domain{N} = \Setabs{a\in \Domain{M}}{\Assign{R}{M}(a, b_1, \dots, b_n)}$
  (veja \olref[bas][sub]{rem:substructure}), e
  $\Expan{M}{X, b_1, \ldots, b_n}$ é a expansão de $\Struct{M}$ que
  interpreta $R$, $c_1$, \dots, $c_n$ por $X$, $b_1,$ \dots, $b_n$,
  respectivamente (com $X \subseteq M^{n+1}$).
\end{enumerate}
\end{defn}
 
\begin{defn}
Dadas duas lógicas abstratas $\tuple{L_1, \models_{L_1}}$ e
$\tuple{L_2, \models_{L_2}}$, dizemos que a segunda é \emph{ao menos
  tão expressiva} quanto a primeira, o que escrevemos
$\tuple{L_1, \models_{L_1}} \leq \tuple{L_2, \models_{L_2}}$, se, para
cada !!{language} $\Lang{L}$ e cada !!{sentence}
$!E \in L_1(\Lang{L})$, existe !!a{sentence}
$!F \in L_2(\Lang{L})$ tal que
$\Mod[L](L_1){!E} = \Mod[L](L_2){!F}$. As lógicas
$\tuple{L_1, \models_{L_1}}$ e $\tuple{L_2, \models_{L_2}}$ são
\emph{equivalentes} se
$\tuple{L_1, \models_{L_1}} \leq \tuple{L_2, \models_{L_2}}$ e
$\tuple{L_2, \models_{L_2}} \leq \tuple{L_1, \models_{L_1}}$.
\end{defn}

\begin{rem}
  A lógica de primeira ordem, isto é, a lógica abstrata
  $\tuple{F, \models}$, é normal. De fato, as propriedades acima são,
  em sua maioria, imediatas para a lógica de primeira ordem. Observamos
  apenas que a propriedade da expansão se reduz à extensionalidade e
  que a relativização de !!a{sentence} $!E$ a
  $\Atom{R}{x, c_1, \dots, c_n}$ é obtida substituindo cada
  !!{subformula} $\lforall[x][!F]$ por
  $\lforall[x][(\Atom{R}{x, c_1, \dots, c_n} \lif \ !F)]$. Além disso,
  se $\tuple{L, \models_L}$ é normal, então
  $\tuple{F, \models} \leq \tuple{L, \models_{L}}$, como se pode
  demonstrar por indução nas !!{formula}s de primeira ordem. Assim, sem
  perda de generalidade, podemos supor que toda !!{sentence} de
  primeira ordem pertence a toda lógica normal.
\end{rem}

\end{document}

